\section{Conclusion and Future Work}

\subsection{Summary of Contributions}

This study presents the first systematic comparison of next-generation YOLO variants (YOLOv8, YOLOv10, YOLOv11, YOLOv12) for real-time durian leaf disease detection on Raspberry Pi edge devices. Through comprehensive evaluation of eight model variants across four optimization formats (PyTorch, ONNX, TFLite, OpenVINO), we provide critical insights into the practical implications of recent YOLO architectural innovations for resource-constrained agricultural applications.

Our key findings demonstrate that:

\begin{enumerate}
\item \textbf{YOLOv11n-ONNX achieves optimal balance}: With 87.1\% mAP@50 and 13.2 FPS on Raspberry Pi 4B, YOLOv11n-ONNX provides the best accuracy-speed trade-off for real-time edge deployment, validating the effectiveness of YOLOv11's refined C3k2 blocks and improved attention mechanisms \cite{4}, \cite{25}.

\item \textbf{ONNX optimization is critical}: ONNX consistently delivers 1.85$\times$ speedup over PyTorch with 99.7\% accuracy retention, highlighting the importance of model optimization for edge deployment. This performance advantage stems from ONNX Runtime's graph optimizations and efficient memory management.

\item \textbf{Nano models enable real-time performance}: All nano models with ONNX/OpenVINO optimization achieve $>$12 FPS, exceeding the 10 FPS threshold for practical real-time deployment. In contrast, small models achieve only 6--7 FPS, limiting their applicability for real-time scenarios on Raspberry Pi 4B.

\item \textbf{TFLite offers storage efficiency}: TFLite INT8 quantization reduces model size by 51.6\% with 2.2\% accuracy loss, making it attractive for storage-constrained devices or bandwidth-limited model distribution, though quantization-aware training could potentially mitigate accuracy degradation \cite{43}.

\item \textbf{Class-specific challenges exist}: Leaf Spot detection (81.3--87.5\% AP@50) remains more challenging than Leaf Blight (88.9--93.1\% AP@50) due to small lesion size and low contrast, suggesting opportunities for targeted architectural improvements or data augmentation strategies.
\end{enumerate}

These findings establish practical benchmarks for deploying YOLO-based disease detection systems in real-world agricultural settings and provide actionable guidelines for practitioners seeking to balance accuracy, speed, and resource constraints in edge AI applications.

\subsection{Broader Implications}

This work demonstrates that state-of-the-art deep learning models can be effectively deployed on affordable (\$35--\$75), resource-constrained edge devices for real-time agricultural monitoring, addressing critical barriers to AI adoption in precision agriculture. The methodologies, benchmarks, and insights presented here extend beyond durian disease detection, providing generalizable guidelines for deploying YOLO-based object detection systems across diverse agricultural applications (pest detection, fruit counting, weed identification, crop health monitoring) and edge computing scenarios (wildlife monitoring, industrial inspection, autonomous vehicles).

Our findings have several important implications:

\begin{itemize}
\item \textbf{Democratizing Agricultural AI}: By demonstrating practical deployment on low-cost hardware, this work reduces financial barriers to AI adoption for smallholder farmers and resource-limited agricultural operations in developing countries.

\item \textbf{Enabling Offline Operation}: Edge-based inference eliminates dependency on stable internet connectivity, making AI-powered disease detection viable in remote rural areas with limited infrastructure.

\item \textbf{Preserving Data Privacy}: On-device inference addresses privacy concerns related to transmitting potentially sensitive farm data to cloud services, an important consideration for commercial agricultural operations.

\item \textbf{Reducing Operational Costs}: Edge deployment eliminates cloud inference costs and reduces bandwidth requirements, improving the economic sustainability of AI-powered agricultural monitoring systems.

\item \textbf{Informing Model Development}: Our systematic evaluation reveals that recent YOLO architectural innovations (C3k2 blocks, improved attention mechanisms) provide measurable benefits for agricultural object detection, guiding future model development efforts.
\end{itemize}

\subsection{Limitations}

While this study provides valuable insights, several limitations should be acknowledged:

\begin{enumerate}
\item \textbf{Dataset Scope}: Our dataset of 2,000 images, while substantial, covers only four disease classes from a limited geographic region (Southern Thailand). Generalization to other durian varieties, disease types, or geographic regions may require additional data collection and model fine-tuning.

\item \textbf{Controlled Evaluation Conditions}: Raspberry Pi benchmarking was conducted under controlled laboratory conditions (stable temperature, no background processes). Real-world field deployment may encounter thermal throttling, power fluctuations, and concurrent system processes that affect performance.

\item \textbf{Single Hardware Platform}: Evaluation focused exclusively on Raspberry Pi 4B. Performance characteristics may differ on alternative edge platforms (Raspberry Pi 5, NVIDIA Jetson Nano, Google Coral, Intel NUC) with different CPU architectures, memory configurations, or specialized AI accelerators.

\item \textbf{Static Image Evaluation}: Our evaluation uses static images rather than continuous video streams. Real-time video processing introduces additional challenges (frame buffering, temporal consistency, motion blur) not addressed in this study.

\item \textbf{Limited Environmental Diversity}: Dataset images were captured under natural lighting conditions but may not fully represent extreme scenarios (heavy rain, fog, dust, direct sunlight, nighttime) encountered in operational agricultural environments.

\item \textbf{No Field Validation}: While our results demonstrate technical feasibility, we have not conducted long-term field trials to validate system reliability, robustness, and practical utility in real agricultural settings.
\end{enumerate}

\subsection{Future Research Directions}

Building on this work, we identify several promising directions for future research:

\subsubsection{Advanced Model Optimization}

\begin{itemize}
\item \textbf{Quantization-Aware Training (QAT)}: Investigate QAT to minimize accuracy degradation from INT8 quantization, potentially improving TFLite performance to match ONNX while maintaining storage efficiency \cite{43}.

\item \textbf{Structured Pruning}: Apply structured pruning and channel pruning to further reduce model complexity while preserving accuracy, potentially achieving 15--20 FPS on Raspberry Pi 4B \cite{2}, \cite{46}.

\item \textbf{Knowledge Distillation}: Explore knowledge distillation from larger teacher models (YOLOv11s, YOLOv11m) to nano student models, potentially improving accuracy without increasing inference cost \cite{19}, \cite{48}.

\item \textbf{Neural Architecture Search (NAS)}: Apply hardware-aware NAS to discover Pareto-optimal architectures specifically optimized for Raspberry Pi's ARM Cortex-A72 CPU \cite{36}, \cite{37}, \cite{38}, \cite{40}.
\end{itemize}

\subsubsection{Hardware Acceleration}

\begin{itemize}
\item \textbf{Raspberry Pi 5 Evaluation}: Benchmark performance on Raspberry Pi 5 with improved CPU (Cortex-A76) and optional AI accelerator, potentially achieving 20--25 FPS with nano models.

\item \textbf{Alternative Edge Platforms}: Evaluate performance on NVIDIA Jetson Nano (GPU-accelerated), Google Coral TPU (specialized AI accelerator), and Intel Neural Compute Stick (VPU-accelerated) for comparative analysis \cite{20}, \cite{47}.

\item \textbf{FPGA Acceleration}: Investigate FPGA-based acceleration for ultra-low-latency applications, potentially achieving $<$10 ms inference time \cite{20}.

\item \textbf{Specialized AI Accelerators}: Explore emerging edge AI accelerators (Hailo-8, Kneron KL720, Rockchip NPU) that promise 10--100$\times$ speedup over CPU inference.
\end{itemize}

\subsubsection{Multi-Modal and Multi-Task Learning}

\begin{itemize}
\item \textbf{Multi-Spectral Imaging}: Integrate visible, near-infrared (NIR), and thermal imaging for improved disease detection under challenging conditions (early-stage symptoms, occlusion, variable lighting) \cite{15}, \cite{63}.

\item \textbf{Multi-Task Learning}: Develop multi-task models jointly performing disease detection, severity assessment, and growth stage classification, potentially improving efficiency and accuracy through shared representations.

\item \textbf{Temporal Modeling}: Incorporate temporal information (LSTM, Transformer) for disease progression prediction and early warning systems, enabling proactive intervention.

\item \textbf{Semi-Supervised Learning}: Leverage unlabeled field data through self-supervised learning and semi-supervised learning to improve model generalization and reduce annotation costs \cite{53}.
\end{itemize}

\subsubsection{System Integration and Field Deployment}

\begin{itemize}
\item \textbf{Complete IoT System}: Develop integrated IoT system combining Raspberry Pi-based detection with cloud dashboard for farm management, automated alerts, and decision support.

\item \textbf{Mobile Application}: Create smartphone application for on-demand disease detection using optimized TFLite models, enabling farmer-friendly diagnostic tools.

\item \textbf{Long-Term Field Trials}: Conduct multi-season field trials to validate system reliability, robustness, and practical utility in real agricultural settings, assessing factors such as weather resistance, power consumption, and maintenance requirements.

\item \textbf{User Acceptance Studies}: Evaluate farmer trust, acceptance, and adoption barriers through participatory design approaches, ensuring developed systems meet practical needs and constraints.

\item \textbf{Edge-Cloud Hybrid Architecture}: Implement hybrid edge-cloud systems where edge devices perform initial screening and cloud services handle complex cases, optimizing accuracy-latency-cost trade-offs \cite{13}, \cite{61}.
\end{itemize}

\subsubsection{Generalization and Transfer Learning}

\begin{itemize}
\item \textbf{Cross-Variety Transfer}: Investigate transfer learning approaches for adapting models to new durian varieties (Monthong, Musang King, D24, Chanee) with minimal additional training data.

\item \textbf{Cross-Crop Transfer}: Explore transfer learning to related tropical fruits (jackfruit, breadfruit, mangosteen) and other crops, developing foundation models for agricultural disease detection.

\item \textbf{Domain Adaptation}: Develop domain adaptation techniques to handle distribution shifts between training and deployment environments (different regions, seasons, camera types, lighting conditions).

\item \textbf{Few-Shot Learning}: Investigate few-shot learning for rapid adaptation to novel disease classes with limited labeled data (5--10 examples), enabling quick response to emerging diseases.
\end{itemize}

\subsubsection{Explainability and Trust}

\begin{itemize}
\item \textbf{Attention Visualization}: Integrate attention visualization and saliency mapping (Grad-CAM, LIME) to provide interpretable disease detection results for agronomists, building trust and facilitating error diagnosis.

\item \textbf{Uncertainty Quantification}: Develop uncertainty quantification methods (Monte Carlo dropout, ensemble models) to flag low-confidence predictions requiring human verification, improving system reliability.

\item \textbf{Active Learning}: Implement active learning frameworks where the system identifies uncertain or novel cases for expert annotation, continuously improving model performance through human-in-the-loop learning.
\end{itemize}

\subsection{Concluding Remarks}

This work establishes a foundation for practical edge AI deployment in precision agriculture, demonstrating that state-of-the-art YOLO models can be effectively deployed on affordable, resource-constrained devices for real-time agricultural monitoring. Our systematic evaluation of YOLOv8--YOLOv12 variants across multiple optimization formats provides actionable guidelines for practitioners and establishes benchmarks for future research.

The recommended configuration (YOLOv11n-ONNX) achieves 87.1\% mAP@50 at 13.2 FPS on Raspberry Pi 4B, representing a practical balance between accuracy, speed, and resource efficiency suitable for real-world deployment. This performance level enables automated disease detection systems that can operate continuously in field conditions, providing early warnings and supporting data-driven decision-making for disease management.

As YOLO architectures continue to evolve and edge hardware capabilities improve, we anticipate further advances in agricultural edge AI. The methodologies and insights presented here provide a foundation for these future developments, contributing to the broader goal of sustainable, efficient, and intelligent agricultural systems that can help address global food security challenges while minimizing environmental impact.

We will release our dataset, trained model weights, and inference code as open-source resources to facilitate reproducibility and accelerate research in agricultural edge AI, fostering collaboration and innovation in this critical application domain.
